\section{Endgame Challenge by Nunn}
\begin{multicols}{2}
    \section*{1}
    \newchessgame[setfen=8/2p5/6K1/8/8/5k2/P7/8]
    \chessboard

    \mythoughts{
        \variation{1. a4 c5} leads to a draw, as both pawns promote.
    }

    So I must play \mainline[level=1]{1. Kf5}.

    \mainline[level=1]{1... Ke3}

    \mythoughts{
        \variation{2. a4} is still a draw, for the same reason as on the previous move. I have no better try.
    }

    \mainline[level=1]{2. Ke5 c6}

    \variation{2... c5 3. Kd5} would be an easy win for White.

    \mainline[level=1]{3. a4 Kd3}

    Both pawns now race to promotion:

    \mainline[level=1]{
        4. a5 c5 5. a6 c4 6. a7 c3 7. a8=Q c2 
    }

    \chessboard

    \mythoughts{
       Queen against a c-pawn is often drawn. If Black's king reaches d2,
       the game is dead drawn, so I must not allow that setup.
    }

    \mainline[level=1]{8. Qd5+ Ke3}
    
    \mainline[level=1]{9. Qg2 Kd3 10. Qg5 Kc3 11. Qc1 Kd3 12. Qb2 Kd2 13. Ke4 Kd1 14. Kd3 c1=Q 15. Qe2# }
    
    \myreflection{
        This is far from a simple position. Up to the queen-versus-pawn phase I
        found the right moves by elimination, but once the queen appeared the winning plan was
        no longer obvious---I did not see the correct continuation.

        A useful reference pattern for such endgames is the fortress with Black's king on d2
        and the pawn still on c2. If Black reaches that setup, the game is usually drawn;
        to play for a win, White must prevent that formation from arising.
    }

    \section*{2}
    \newchessgame[setfen=8/8/1KP5/3r4/8/8/8/k7 w - - 0 1]
    \chessboard

    \mythoughts{
        This is a rare motif: a lone pawn outwits a rook. The first
        move is obvious.
    }

    \mainline[level=1]{1. c7 Rd6+}

    \mythoughts{
        I cannot play \variation{2. Kc5 Rd2}: Black would have perpetual
        checks on the a- and b-files. To escape the checks, White would need to bring the king
        to c3 (otherwise \symrook d2+ wins the c2-pawn). Black could then tuck the king on b1
        and follow up with \symrook c2, picking off the pawn.

        White must therefore march the king down the board while Black keeps checking from behind.
    }

    \mainline[level=1]{
2. Kb5 Rd5+ 3. Kb4 Rd4+ 4. Kb3 Rd3+ 5. Kc2 Rd4 
    }

    The remainder is study-like:
    \mainline[level=1]{
        6. c8=R Rd2+ 7. Kxd2 Kb2 8. Rc3 Kb1 9. Rb3+ Ka1 10. Kc2 Ka2 11. Rd3 Ka1 12. Ra3# 
    }

    \myreflection{
        Black is not obliged to follow this line; for instance, Black could simply allow White
        to promote to a queen. In over-the-board play I would be unlikely to convert
        queen versus rook from scratch, so more study is needed.
    }

    \section*{3}
    \newchessgame[setfen=5K2/2p5/8/2k5/8/2P5/4P3/8 w - - 0 1]
    \chessboard

    \mythoughts{
        Black has a c-pawn, so a queen-versus-pawn endgame is likely to arise.
        I do not yet have a clear sense of when such an endgame is winning and when it is not.

        Let me start by eliminating some moves.

        \variation{1. e4 Kd6}.

        \variation{1. Ke7 Kc4 2. Kd7 c5 3. e4 Kxc3 4. e5} Black's king
        can go to either b2 or d2. White has three more moves before the pawn queens,
        and Black can push to c2 right after White promotes. The position is a draw.

        \variation{
            1. Kf7 Kc4? 2. e4 Kxc3 3. e5 c5 4. e6 c4 5. e7 Kd2 6. e8=Q c3 7. Qd8+ Kc1 8. Qg5+ Kb1 
        }
        A queen generally wins against a pawn on the 6th rank, so no further calculation is needed.

        \variation{
            1. Kf7 Kd5 2. Kf6 Kc4 3. e4 Kxc3 4. e5 c5 5. e6 Kd2 6. e7 c4 7. e8=Q c3 8. Qd8+ Kc1 
        }

    } 

    \myreflection{
        In the two \variation{1. Kf7} lines, it is instructive to compare the resulting positions:
        one is winning, the other is not. The only difference is whether White's own king
        blocks the queen from giving the decisive check on g5.
    }

    \chessboard[
        setfen=3Q4/5K2/8/8/8/2p5/8/2k5 w - - 2 8
    ]

    \chessboard[
        setfen=3Q4/8/5K2/8/8/2p5/8/2k5 w - - 2 9
    ]

    This means the only correct move is \mainline[level=1]{1. Kg7}.

    \variation[invar]{
        1... Kc4
    } simply loses to \variation{2. e4}. We reach a similar position
    when White's pawn queens with Black's pawn still on the 5th rank: White can give
    the check on g5 and is therefore winning.

    \mainline[level=1]{
        1... Kd5
    }

    \chessboard

    \mythoughts{
        Pawn moves are meaningless here. \variation[invar]{2. Kf6} transposes to a position
        we saw earlier, where White cannot give the check on g5. That leaves only one move.
    }

    \mainline[level=1]{
        2. Kf7
    }

    \myreflection{
        The following position is a reciprocal zugzwang: since Black is to move,
        Black loses.
    }

    \chessboard[
        setfen=8/2p2K2/8/3k4/8/2P5/4P3/8 b - - 3 2
    ]

    \variation[invar]{
        2... Kc4
    } again loses, because once White queens he can give the familiar check on g5.

    \mainline[level=1]{2... Kd5 3. Ke7 Ke5 4. Kd7 Kd5 5. Kc7 c4 6. Kd7 Ke5 7. Kc6
    }

    White will pick up the pawn and win.

    \myreflection{
        This is certainly a difficult problem. The queen-versus-pawn endgame is still
        not clear to me; I need to study it further.
    }

    \section*{4}
    \newchessgame[
        setfen=3K4/kp6/p7/1P6/8/8/7P/8 w - - 0 1
    ]
    \chessboard

    \mythoughts{
        \variation{1. bxa6 b5} both pawns queen. The position is a draw.

        The first move is obvious.
    }

    \mainline[level=1]{
        1. b6+ 
    }

    Black cannot take the pawn because when White promotes, White's queen
    also controls the promotion square for the a pawn. The resulting endgame is loss 
    for Black.

    \mainline[level=1]{
        1... Kb8
    }

    The next moves are obvious.
    \mainline[level=1]{
        2. h4 a5 3. h5 a4 4. h6 a3 5. h7 a2 6. h8=Q a1=Q!
    }

    \chessboard

    \mythoughts{
        Taking the queen is stalemate. 
    }

    Here I fail to find the right move. White's queen should stay 
    on the 8th rank. Threatening a checkmate by a king move. Black cannot 
    prevent it with a check. The only move is therefore

    \mainline[level=1]{
        7. Qg8 Qa2
    }

    \mythoughts{
       \variation[invar]{8. Qf8 Qa3}
    } fails. Moving the queen away allows Black to check on d6. Taking the queen is again stalemate.
    
    White's queen must still stay on the 8th rank. The only move is therefore

    \mainline[level=1]{
        8. Qe8 Qa4
    }

    \chessboard

    Going back to h8 doesn't make any progress because Black can repeat by moving 
    the queen to a1.

    \mainline[level=1]{
        9. Qe5+ Ka8 10. Qh8
    } 

    Black is in zugzwang and therefore loses:

    \mainline[level=1]{
        10... Qf4 11. Qa1 Kb8 12. Qa7#
    } 

\end{multicols}     